\section{Constructive Neural Bellman Operators beyond quadratic continuations}
\label{sec:r39nonlinear}

A continuation method has an economic use when its approximation can be
translated into a decision guarantee without assuming that the approximation
class contains the value function. We now give this translation for nonlinear
recursive economies. The comparison is with all feasible adapted policies.
The construction uses compact state and action sets, not a Gaussian coefficient
identity. The preceding diffusion framework and unbounded-state results remain
in force under their respective assumptions; compactness is imposed only in
this additional realization.

\subsection{A residual-oscillation theorem}
Let $X$ and $A$ be compact metric spaces, let $F_t:X\times A\times Z\to X$
and $c_t:X\times A\to\R$ be continuous, and let $Z$ have a finite distribution
with fixed probabilities $q_z$. Write $\beta_t>0$ for the date-$t$ continuation
weight, $B_0=1$, and $B_t=\prod_{j<t}\beta_j$. The conditional continuation
functional $\rho_t$ is monotone and cash invariant:
$\rho_t(Y+a)=\rho_t(Y)+a$. It is continuous on finite vectors. In particular,
these conditions hold for
\begin{equation}\label{eq:r39risk}
 \rho_\theta(Y)=\theta^{-1}\log\sum_zq_z e^{\theta Y(z)},\qquad
 \rho_0(Y)=\sum_zq_zY(z),\qquad \theta\geq0.
\end{equation}
Boundedness makes the exponential moment finite. Define
\[
 Q_t(v;x,a)=c_t(x,a)+\beta_t\rho_t(v(F_t(x,a,Z))),\qquad
 T_tv(x)=\min_{a\in A}Q_t(v;x,a).
\]
Continuity gives measurable minimizers. Recursive costs for arbitrary adapted
policies use the same conditional functional at every realized history. Let
$V_T^*=g$ and $V_t^*=T_tV_{t+1}^*$.

\begin{theorem}[Nonlinear full-policy comparison]\label{thm:r39comparison}
Let $v_t$ be continuous functions and let $\pi_t(x)\in A$ be measurable.
Suppose that, at every state,
\begin{align}
 \ell_t\leq v_t-T_tv_{t+1}\leq u_t,&\qquad t<T,\label{eq:r39residual}\\
 Q_t(v_{t+1};x,\pi_t(x))-T_tv_{t+1}(x)&\leq\eta_t,\label{eq:r39actor}\\
 \ell_T\leq v_T-g&\leq u_T.\nonumber
\end{align}
Then, for every initial state and against all feasible adapted policies,
\begin{equation}\label{eq:r39loss}
 0\leq J_0^\pi-V_0^*
 \leq\sum_{t<T} B_t(u_t-\ell_t+\eta_t)+B_T(u_T-\ell_T).
\end{equation}
Moreover, $v_0-\sum_{t\leq T}B_tu_t$ is a lower bound on $V_0^*$.
Neither membership of $V^*$ in a neural class nor convergence of a training
optimizer is required for this comparison.
\end{theorem}

\begin{proof}
Monotonicity and cash invariance imply
$\|T_tv-T_tw\|_\infty\leq\beta_t\|v-w\|_\infty$; the same statement holds
for the operator of any fixed feasible action. Backward induction applied to
\eqref{eq:r39residual} yields
\[
 v_t-\sum_{s=t}^T(B_s/B_t)u_s\leq V_t^*
 \leq v_t-\sum_{s=t}^T(B_s/B_t)\ell_s.
\]
The terminal equality is immediate. For the induction step, shift the next
function by the corresponding deterministic constant, apply $T_t$, and use
$T_tv_{t+1}\in[v_t-u_t,v_t-\ell_t]$.
For the returned policy, \eqref{eq:r39actor} instead gives
\[
 J_t^\pi\leq v_t+\sum_{s=t}^{T-1}(B_s/B_t)(\eta_s-\ell_s)
                    -(B_T/B_t)\ell_T.
\]
Subtract the lower bound on $V_t^*$. At an arbitrary history, each realized
action has continuation cost at least the Bellman minimum, and induction
compares its future cost with $V_{t+1}^*$. Thus the Bellman value is a lower
bound for every adapted policy, not only for the Markov or neural class.
A measurable minimizing Markov policy attains it. This proves both claims.
\end{proof}

The economically relevant residual is its oscillation $u_t-\ell_t$.
A common additive valuation error cancels from the loss comparison. This is
the nonlinear counterpart of centered continuation error in the earlier
operator analysis, rather than an assumption that a critic learns absolute
value levels accurately.

\subsection{A finite construction and a trained-network check}
Take $X=[0,1]$ and $A=[0,\bar a]$. Suppose $c_t$ and $F_t$ have known
Lipschitz constants separately in state and action. If $v_{t+1}$ has
Lipschitz constant $L_{t+1}$, admissible constants for $Q_t$ are
\begin{equation}\label{eq:r39lipschitz}
 L_{x,t}=L_{c,x,t}+\beta_t L_{F,x,t}L_{t+1},\qquad
 L_{a,t}=L_{c,a,t}+\beta_t L_{F,a,t}L_{t+1}.
\end{equation}
The same $L_{x,t}$ bounds the state variation of $T_tv_{t+1}$.
Use a uniform state partition $x_i=ih$ and an action grid of mesh $\Delta$,
including both endpoints. At every pair, verified arithmetic supplies
$\underline q_{ij}\leq Q_t(v_{t+1};x_i,a_j)\leq\overline q_{ij}$.
Set $L_i=\min_j\underline q_{ij}$, $U_i=\min_j\overline q_{ij}$ and
$\delta_t=L_{a,t}\Delta/2$. Let $v_t$ be the linear interpolant of stored
nodal values $y_i$, with no requirement that these equal an exact minimizer.
Then valid residual constants are
\begin{align}\label{eq:r39mesh}
 \ell_t&=\min_i(y_i-U_i)-L_{x,t}h/2,\\
 u_t&=\max_i(y_i-L_i+\delta_t)+L_{x,t}h/2.\nonumber
\end{align}
If the action $a_{j_i}$ is stored at node $i$ and extended by the
nearest-node rule, a valid actor allowance is
\begin{equation}\label{eq:r39meshactor}
 \eta_t=\max_i(\overline q_{i j_i}-L_i)+\delta_t+L_{x,t}h.
\end{equation}
All expressions are enclosed outward in the implementation.

\begin{proposition}[Constructive nonquadratic continuation]\label{prop:r39construction}
Under the preceding conditions, \eqref{eq:r39mesh}--\eqref{eq:r39meshactor}
and a certified terminal interpolation error imply
Theorem~\ref{thm:r39comparison}. Every such $v_t$ has the exact representation
\begin{equation}\label{eq:r39relu}
 v_t(x)=y_0+s_0x+\sum_{i=1}^{N-1}(s_i-s_{i-1})(x-x_i)_+,
 \qquad s_i=(y_{i+1}-y_i)/h.
\end{equation}
For each fixed finite horizon, a feasible neural policy with any prescribed
positive loss tolerance can be constructed in finitely many operations if
arithmetic precision can be increased to deliver the required interval widths.
This assertion concerns the constructive mesh algorithm, not unrestricted
stochastic-gradient training.
\end{proposition}

\begin{proof}
The continuous action minimum at a node lies in $[L_i-\delta_t,U_i]$:
a grid point is within $\Delta/2$ of any minimizing action. On a state cell,
linear interpolation of nodal Bellman values differs from the Bellman function
by at most $L_{x,t}h/2$, since the interpolation-weighted distance to the two
endpoints is $2r(h-r)/h\leq h/2$. This proves \eqref{eq:r39mesh}.
Moving from a state to its nearest node changes both its chosen-action value
and the Bellman minimum by at most $L_{x,t}h/2$ each, proving
\eqref{eq:r39meshactor}. Slope changes prove \eqref{eq:r39relu}.

For finiteness, choose nodal midpoints of $[L_i,U_i]$ and the action minimizing
$\overline q_{ij}$. Suppose each $Q$ interval has width at most $\omega$ and
each stored midpoint has an additional rounding error at most $\zeta$.
Then $U_i-L_i\leq\omega$. Residual oscillation plus actor allowance is at most
\[
 2\omega+2\zeta+2\delta_t+2L_{x,t}h.
\]
For the Lipschitz recursion, each nodal midpoint is within
$\delta_t+\omega+\zeta$ of the exact Bellman value. Consequently the next
interpolant has Lipschitz constant at most
$L_{x,t}+2(\delta_t+\omega+\zeta)/h$.
Choose $\Delta\leq h^2$ and $\omega+\zeta\leq h^2$ at every date.
Equations~\eqref{eq:r39lipschitz} then give, by finite backward induction,
a bound on all interpolant Lipschitz constants uniform over $0<h\leq1$.
Terminal interpolation contributes at most $L_g h/2+\zeta$ in absolute error.
The right-hand side of \eqref{eq:r39loss} therefore tends to zero with $h$.
A sufficiently fine finite grid and sufficiently precise finite arithmetic
attain any positive target. This proof does not assert that fixed binary64
arithmetic suffices for every target.
\end{proof}

A conventional spline routine and \eqref{eq:r39relu} represent the same
constructed function; relabeling it alone is not a neural work advantage.
The second implementation trains all affine hidden weights in
$f_t(x)=a_t+b_tx+\sum_{j=1}^p c_{tj}(w_{tj}x+d_{tj})_+$.
It uses the certified nodal values as deterministic training data and selects
actions from its own continuation predictions. Its selected actions are then
checked by \eqref{eq:r39meshactor} against the independent interval brackets.
The full cost of the reference construction, training and verification is
charged. Acceptance certifies the returned policy even when the optimizer has
no convergence guarantee. A rejected proposal can be replaced by the
constructive policy above; this gives a terminating certified procedure without
asserting that every trained proposal must be accepted.

\begin{proposition}[All-state verification of trained ReLU continuations]
\label{prop:r39breakpoints}
Let $f$ be the preceding one-hidden-layer scalar ReLU network and $v$ a linear
spline. The maximum of $|f-v|$ on $[0,1]$ occurs at an endpoint, a spline knot,
or a neural root $-d_j/w_j$ in that interval, with $w_j\ne0$. Enclosing each
root by outward division and evaluating both functions over its interval gives
a certified uniform bound, including nondyadic roots. Direct evaluation costs
$O((N+p)(p+\log N))$ arithmetic operations; sorting the breakpoints, when used,
costs $O((N+p)\log(N+p))$.
\end{proposition}
\begin{proof}
Between consecutive points in the specified set, $f-v$ is affine. The absolute
value of an affine function is convex and is bounded by its endpoint maximum.
Continuity includes the endpoints of adjacent pieces. Zero hidden slopes
produce no breakpoint. Every exact breakpoint belongs to its outward root
interval, so interval evaluation bounds its value even when a root cannot be
stored exactly. A direct network evaluation uses $O(p)$ operations and a
binary-search spline evaluation $O(\log N)$. There are at most $N+p+1$ points.
\end{proof}

This certificate removes the mesh-Lipschitz remainder used in the earlier
network-to-spline check. It does not retroactively change the earlier policy
bounds or execution clocks. The new endpoint audit uses exact rational
arithmetic only as an independent check of the verifier, not as the production
method used in the floating-point economic services.

\subsection{Nonlinear investment and the value of reoptimization}
\label{sec:r39economy}
The economic realization has bounded capital, irreversible investment and a
shortfall penalty. Its primitives are
\begin{align}
 F(x,a,z)&=\operatorname{clip}_{[0,1]}
       \{13x/16+a+x(1-x)/16+z\},\label{eq:r39capital}\\
 c_P(x,a)&=(x-11/16)^2+Pa^2+4a^4+2(3/8-x)_+^2,\nonumber\\
 g(x)&=2(x-11/16)^2+2(3/8-x)_+^2,\qquad \beta=15/16,\nonumber
\end{align}
where $a\in[0,1/4]$ and the shocks $(-1/16,0,1/16)$ have probabilities
$(1/4,1/2,1/4)$. The price $P$ measures the quadratic investment burden;
$\theta$ in \eqref{eq:r39risk} measures aversion to high future costs.
A low-capital household faces both a distance-to-target loss and a shortfall
penalty. Investment is capped and cannot be negative. The quartic adjustment
cost, nonlinear accumulation and clipping prevent a Riccati reduction.
This is a theoretical counterfactual economy, not an estimated calibration.

Valid primitive constants are $L_{c,x}=23/8$, $L_{F,x}=7/8$,
$L_{c,a}=P/2+1/4$, $L_{F,a}=1$, and $L_g=17/4$.
They follow by differentiation on smooth pieces and continuity at the kinks;
clipping is nonexpansive. Hence Proposition~\ref{prop:r39construction}
verifies its assumptions directly from the primitives.

\begin{proposition}[Fixed-rule investment exposure]\label{prop:r39exposure}
For a fixed feasible policy rule $\pi$ in \eqref{eq:r39capital}, independent of
$P$, the derivative $D_t^\pi=\partial J_t^\pi/\partial P$ exists along each
finite shock tree and satisfies $D_T^\pi=0$ and
\begin{equation}\label{eq:r39exposure}
 D_t^\pi=a_t^2+\beta\sum_z\widetilde q_{t,z}D_{t+1}^\pi(F_t),\qquad
 \widetilde q_{t,z}=
 \frac{q_z\exp(\theta J_{t+1}^\pi(F_t))}
      {\sum_yq_y\exp(\theta J_{t+1}^\pi(F_t))}.
\end{equation}
At $\theta=0$, $\widetilde q=q$. In particular $D_t^\pi\geq0$.
For old and new rules, the total cost change decomposes exactly as
\begin{align}\label{eq:r39decomposition}
 J_0^{\pi^{\rm new}}(4)-J_0^{\pi^{\rm old}}(1)
 ={}&\{J_0^{\pi^{\rm old}}(4)-J_0^{\pi^{\rm old}}(1)\}\\
 &-\{J_0^{\pi^{\rm old}}(4)-J_0^{\pi^{\rm new}}(4)\}.\nonumber
\end{align}
The first bracket is mechanical repricing; the second is the cost saved by
reoptimization at the new price.
\end{proposition}
\begin{proof}
For a fixed rule, the states and selected actions on the finite shock tree do
not depend on $P$, even if the rule is discontinuous in state. Each stage cost
is affine in $P$. Differentiate the finite log-sum-exp recursion; its derivative
is the probability-weighted derivative in \eqref{eq:r39exposure}. Backward
induction proves existence and nonnegativity. Add and subtract
$J_0^{\pi^{\rm old}}(4)$ to obtain \eqref{eq:r39decomposition}.
\end{proof}

No global monotonicity of the optimal investment rule is inferred from this
fixed-rule statement. The numerical counterfactual below instead reports the
actual policy changes and an independent lower bound on the second bracket.
A strictly positive lower bound establishes an economic gain from the returned
policy even when its optimality certificate is conservative.
